\documentclass[../../main.tex]{subfiles}
\begin{document}
% SOURCE p=250; printed=235; object=ch08-preliminaries-title; region=middle-lower
\section{准备知识}\label{section:ch08-preliminaries}
% SOURCE p=250; printed=235; object=ch08-function-space-conventions; region=lower
\begin{definition}[本章的函数空间与欧氏记号]\label{definition:ch08-function-spaces}
本章在 $\mathbb R^n$ 上开展讨论，$n$ 始终表示维数。涉及测度时，除非另行说明，均指 Lebesgue 测度。因此，对 $\mathbb R^n$ 中可测集 $E$，将 $L^p(E,m)$ 简写为 $L^p(E)$。若 $U\subseteq\mathbb R^n$ 开、$k\in\mathbb N$，则 $C^k(U)$ 为 $U$ 上所有不超过 $k$ 阶偏导数均存在且连续的函数组成的空间，且
\begin{align*}
C^\infty(U)=\bigcap_{k=1}^{\infty}C^k(U).
\end{align*}
对任意 $E\subseteq\mathbb R^n$，$C_c^\infty(E)$ 表示所有在\emph{整个 $\mathbb R^n$ 上}属于 $C^\infty$、支撑紧且包含于 $E$ 的函数组成的空间。若 $E=\mathbb R^n$ 或 $U=\mathbb R^n$，通常省略这个定义域：例如 $L^p=L^p(\mathbb R^n)$、$C^k=C^k(\mathbb R^n)$、$C_c^\infty=C_c^\infty(\mathbb R^n)$。对 $x,y\in\mathbb R^n$，置
\begin{align*}
x\cdot y=\sum_{j=1}^n x_jy_j,\qquad |x|=\sqrt{x\cdot x}.
\end{align*}
\end{definition}
% SUPPLEMENT: id=v14-ch08-function-spaces-details; parent=ch08-function-space-conventions; kind=motivation+hypotheses; origin=independent-derivation
\begin{remark}
\textbf{动机与条件（重构新增）。}同时使用可积性、连续性和任意阶微分，是为了随后把积分运算与求导连接起来。$U$ 开使每点附近均可在各坐标方向求通常的双侧导数；$k$ 阶导数连续保证相应阶数内混合偏导的交换，并可用微积分基本定理。$C^0$ 如需使用即指连续函数。$C_c^\infty(E)$ 中的 $E$ 无需开，因为函数的光滑性是在环境空间中要求的；“支撑包含于 $E$”比“在 $E$ 上光滑”更强。若 $E$ 内部为空，则该空间只有零函数：非零连续函数在某个非空开集上不为零，迫使这个开集包含于其支撑及 $E$。紧支撑则保证函数及每阶导数均有界、可积。$L^p$ 的元素仍按几乎处处相等识别，$C^k$ 中讨论的则是逐点定义的函数；这一区别在平移问题中会用到。内积取实坐标的欧氏内积，$n$ 为固定正整数。
\end{remark}
% END SUPPLEMENT: id=v14-ch08-function-spaces-details
% SOURCE p=251; printed=236; object=ch08-multiindex; region=upper
\begin{definition}[多重指标]\label{definition:ch08-multiindex}
记 $\partial_j=\partial/\partial x_j$。\textbf{多重指标}是非负整数的有序 $n$ 元组 $\alpha=(\alpha_1,\ldots,\alpha_n)$。定义
\begin{align*}
|\alpha|=\sum_{j=1}^n\alpha_j,\qquad
\alpha!=\prod_{j=1}^n\alpha_j!,\qquad
\partial^\alpha=\left(\frac{\partial}{\partial x_1}\right)^{\alpha_1}\cdots
\left(\frac{\partial}{\partial x_n}\right)^{\alpha_n},\qquad
x^\alpha=\prod_{j=1}^n x_j^{\alpha_j}.
\end{align*}
这里 $|\alpha|$ 表示各分量之和，与 $|x|=(\sum_jx_j^2)^{1/2}$ 的欧氏长度约定不同；具体含义由上下文确定。
\end{definition}
% SUPPLEMENT: id=v14-ch08-multiindex-details; parent=ch08-multiindex; kind=motivation+hypotheses; origin=independent-derivation
\begin{remark}
\textbf{动机与条件（重构新增）。}多重指标记录每个坐标分别求导多少次，把总阶数相同的项归为一组。分量非负且为整数，是因为这里讨论普通整数阶导数和多项式；$\alpha!$ 是各坐标阶乘的乘积，不是 $|\alpha|!$。约定零次幂、零次导数及 $0!$ 均为单位，因此 $x^0=1$、$\partial^0f=f$。以后 $\beta\leqslant\alpha$ 指逐坐标不等式，$\alpha+\beta$ 也逐坐标相加。这些约定使零阶和常数情形无需排除。
\end{remark}
% END SUPPLEMENT: id=v14-ch08-multiindex-details
% SOURCE p=251; printed=236; object=ch08-taylor-formula; region=middle
\begin{theorem}[多重指标形式的 Taylor 公式]\label{theorem:ch08-taylor-formula}
若 $f\in C^k$，则在任意 $x_0\in\mathbb R^n$ 处有
\begin{align*}
f(x)&=\sum_{|\alpha|\leqslant k}(\partial^\alpha f)(x_0)
\frac{(x-x_0)^\alpha}{\alpha!}+R_k(x),\qquad
\lim_{x\to x_0}\frac{|R_k(x)|}{|x-x_0|^k}=0.
\end{align*}
\end{theorem}
% SUPPLEMENT: id=v14-ch08-taylor-proof; parent=ch08-taylor-formula; kind=omitted-proof+motivation+hypotheses; origin=independent-derivation
\begin{proof}
\textbf{完整补写（重构新增）。}置 $h=x-x_0$。$k=0$ 时结论就是连续性。设 $k\geqslant1$，对 $0\leqslant t\leqslant1$ 置 $g(t)=f(x_0+th)$。重复链式法则给出
\begin{align*}
g^{(j)}(t)=\sum_{|\alpha|=j}\frac{j!}{\alpha!}h^\alpha
(\partial^\alpha f)(x_0+th),\qquad 0\leqslant j\leqslant k.
\end{align*}
系数 $j!/\alpha!$ 计数 $j$ 次求导时分别选中各坐标 $\alpha_1,\ldots,\alpha_n$ 次的排列数，故这一步不预设待证 Taylor 公式。对 $g$ 反复使用微积分基本定理，得到
\begin{align*}
g(1)=\sum_{j=0}^{k-1}\frac{g^{(j)}(0)}{j!}
+\frac1{(k-1)!}\int_0^1(1-t)^{k-1}g^{(k)}(t)\,\mathrm dt.
\end{align*}
这个积分余项可直接归纳验证：$k=1$ 是基本定理；再将最高阶导数写成其零点值加下一阶导数的积分，并在紧三角形上交换连续函数的积分次序，就得到下一阶的核。由于 $\int_0^1(1-t)^{k-1}\,\mathrm dt=1/k$，减去 $g^{(k)}(0)/k!$ 后，余项为
\begin{align*}
R_k(x)=\frac1{(k-1)!}\int_0^1(1-t)^{k-1}
\bigl[g^{(k)}(t)-g^{(k)}(0)\bigr]\,\mathrm dt.
\end{align*}
每个 $|h^\alpha|\leqslant|h|^k$，而有限多个 $k$ 阶偏导数在 $x_0$ 连续。因此
\begin{align*}
|R_k(x)|\leqslant |h|^k
\sum_{|\alpha|=k}\frac1{\alpha!}
\sup_{0\leqslant t\leqslant1}
|\partial^\alpha f(x_0+th)-\partial^\alpha f(x_0)|
=o(|h|^k).
\end{align*}
将各 $g^{(j)}(0)$ 代入便得公式。复值函数也可分别对实部、虚部应用上述证明。
\end{proof}
\begin{remark}
\textbf{动机与条件（重构新增）。}把多变量函数限制到一条线段上，就把局部近似归结为一元微积分；多重指标负责收集同阶项。这里只要求 $C^k$，不需要 $k+1$ 阶导数，因为结论是 $o(|h|^k)$ 而非带 $k+1$ 阶导数的余项估计。连续性使线段上的最高阶导数振幅趋零，正是最后一步的根据。公式也是局部的：若仅 $f\in C^k(U)$ 且 $x_0\in U$，先取以 $x_0$ 为中心且包含于 $U$ 的小球，足够短的线段便仍在 $U$ 内，同样证明成立。
\end{remark}
% END SUPPLEMENT: id=v14-ch08-taylor-proof
% SOURCE p=251; printed=236; object=ch08-leibniz-formula; region=middle
\begin{proposition}[偏导数的乘积法则]\label{proposition:ch08-leibniz-formula}
对有相应阶连续偏导数的 $f,g$，
\begin{align*}
\partial^\alpha(fg)=\sum_{\beta+\gamma=\alpha}
\frac{\alpha!}{\beta!\gamma!}(\partial^\beta f)(\partial^\gamma g).
\end{align*}
\end{proposition}
其完整证明及与单项式相乘的推论见\cref{exercise:ch08-leibniz}。
% SUPPLEMENT: id=v14-ch08-leibniz-motivation; parent=ch08-leibniz-formula; kind=motivation+hypotheses; origin=independent-derivation
\begin{remark}
\textbf{动机与条件（重构新增）。}每次求导只能落在两个因子之一，系数记录这些分配方式。固定 $\alpha$ 时，有不超过 $|\alpha|$ 阶的连续偏导数已经足够；无需为这一个公式假设无限光滑。连续混合偏导使求导次序可以整理成多重指标形式；$\beta+\gamma=\alpha$ 则保证每个坐标的求导次数守恒。$\alpha=0$ 时公式退化为 $fg=fg$。
\end{remark}
% END SUPPLEMENT: id=v14-ch08-leibniz-motivation
% SOURCE p=251; printed=236; object=ch08-function-value-notation; region=middle-lower
我们常采用一种不甚严格但很方便的写法：以同一符号表示函数及它在某点的值。例如，“$x^\alpha$”也可以表示在点 $x$ 取值 $x^\alpha$ 的函数。
% SOURCE p=251; printed=236; object=ch08-bump-function; number=8.1; region=lower
\begin{example}[非零光滑紧支撑函数]\label{example:ch08-bump-function}
$\mathbb R^n$ 上有两个 $C^\infty$ 函数空间对我们尤为重要。第一个是紧支撑光滑函数空间 $C_c^\infty$。其中存在非零函数并非一目了然；标准构造利用
$\eta(t)=e^{-1/t}$（$t>0$）、$\eta(t)=0$（$t\leqslant0$）即使在原点也光滑这一事实，见\cref{exercise:ch08-flat-function}。令
\begin{align}\label{eq:ch08-bump-function}
\psi(x)=\eta(1-|x|^2)=
\begin{cases}
\exp\bigl[(|x|^2-1)^{-1}\bigr],&|x|<1,\\
0,&|x|\geqslant1.
\end{cases}
\end{align}
则 $\psi\in C^\infty$，且 $\operatorname{supp}\psi$ 是闭单位球。下一节将用这个单独的函数构造丰富的 $C_c^\infty$ 元素；见\cref{proposition:ch08-smooth-density,proposition:ch08-smooth-cutoff}。
\end{example}
% SUPPLEMENT: id=v14-ch08-bump-details; parent=ch08-bump-function; kind=omitted-proof+motivation+hypotheses; origin=independent-derivation
\begin{remark}
\textbf{补证与动机（重构新增）。}\cref{exercise:ch08-flat-function}证明 $\eta$ 在零点的所有阶导数为零。因此它与光滑函数 $x\mapsto1-|x|^2$ 的复合在球面上也光滑。$\psi$ 在开单位球严格为正，在球外为零，所以其非零集的闭包恰为闭单位球；该球紧，且 $\psi(0)=e^{-1}>0$。关键是零点处的无限阶平坦性，它使内外两段光滑接合。单纯把一个在球面处导数不消失的函数在球外补零，通常不能做到这一点。
\end{remark}
% END SUPPLEMENT: id=v14-ch08-bump-details
\subfile{section-01-schwartz.tex}
\end{document}
